\subsection{Tensorized Logarithmic-Sobolev Inequality}
\label{app:lsi-proof}

The coordinate estimate in Lemma~\ref{lem:deficit} uses
\[
             \Ent(w^2)\le2\E[w\Lapl w].
\]
This classical inequality measures how far the nonnegative field
\(w^2\) is from constant by the squared differences of \(w\) across
edges. We include its proof, including the step that makes its constant
independent of dimension.

Start with one bit. Replacing the two values of \(w\) by their absolute
values preserves entropy and reduces their squared difference.
Homogeneity then lets us assume \(\E w^2=1\) and write the values as
\(\sqrt{1+s},\sqrt{1-s}\), where \(0\le s\le1\). In this notation
\begin{align*}
 \Ent(w^2)&=\phi(s),\\
 2\E[w\Lapl w]&=\frac{(w_+-w_-)^2}{2}=1-\sqrt{1-s^2}.
\end{align*}
The derivative of \(s/\sqrt{1-s^2}-\atanh s\) is
\((1-s^2)^{-3/2}-(1-s^2)^{-1}\ge0\). The difference vanishes at
zero. Integrating \(\phi'(s)=\atanh s\) therefore proves
\(\phi(s)\le1-\sqrt{1-s^2}\). Continuity supplies \(s=1\).

It remains to tensorize the one-bit bounds. For a finite probability
space, entropy is convex as a function of a nonnegative vector \(z\).
For positive \(z\), its second derivative in a direction \(v\) is
\[
 \E\frac{v^2}{z}-\frac{(\E v)^2}{\E z}\ge0
\]
by Cauchy--Schwarz. Zero entries follow by continuity.
For independent coordinates \(X,Y\), the entropy chain rule and this
convexity give
\begin{align*}
 \Ent_{X,Y}(z)
 &=\E_Y\Ent_X(z)+\Ent_Y(\E_Xz)\\
 &\le\E_Y\Ent_X(z)+\E_X\Ent_Y(z).
\end{align*}
Induction over the coordinates yields
\(\Ent(z)\le\sum_i\E\Ent_i(z)\), where \(\Ent_i\) takes entropy
over coordinate \(i\) with all others fixed. Apply the one-bit result
on each such pair, with \(z=w^2\):
\[
 \Ent(w^2)\le\frac12\sum_i\E_{-i}(w_{i,+}-w_{i,-})^2
                         =2\E[w\Lapl w].
\]
No constant was multiplied by the number of coordinates.

\subsection{Propagation of the Local Cutoff}
\label{app:local-propagation}
This is the threshold-propagation property in
Yu's Lemma 4.9~\cite{Yu}, expressed in sign means and natural logarithms.
We include the one-sign-change proof so that the local comparison
requires no numerical monotonicity test or external lemma.

Define the scalar local margin
\[
 M_{\rm loc}(r,a)=(1-r^2)h(ar)-(1-ar^2)h(r),\qquad 0\le r,a\le1.
\]
We prove the propagation used below: if \(M_{\rm loc}(r_1,a)\ge0\) for some
\(0<r_1<1\), then \(M_{\rm loc}(r,a)\ge0\) whenever \(0\le r\le r_1\).
The reason is a single change of sign in its entropy-series coefficients.

The case \(a=1\) is an identity. For \(a<1\), put \(t=r^2\) and
\(H(t)=h(\sqrt t)\). Let \(b_n=1/[2n(2n-1)]\) be the entropy-series
coefficient, and let \(T_n=\sum_{j>n}b_j\) be its tail. Then
\begin{align*}
 H(t)&=L-\sum_{n\ge1}b_nt^n,\\
 \frac{H(t)}{1-t}&=\sum_{n\ge0}T_nt^n,\\
 T_n&=\int_0^1\frac{x^{2n}}{1+x}\,\mathrm dx.
\end{align*}
For the last identity, express
\(b_j=\int_0^1x^{2j-2}(1-x)\,\mathrm dx\), sum the geometric series,
and cancel \(1-x\). All series used here converge absolutely for
\(t<1\).

Expansion of the local margin now gives
\begin{align*}
 \frac{M_{\rm loc}(\sqrt t,a)}{1-t}&=\sum_{n\ge1}q_nt^n,\\
 q_n&=\frac{1-a}{2n}\left[
 \frac{a+a^2+\cdots+a^{2n-1}}{2n-1}-2nT_n\right].
\end{align*}
The average of the powers of \(a\) decreases with \(n\): the next
average appends two terms no larger than the preceding terms.
In the opposite direction,
\[
 (n+1)T_{n+1}-nT_n
 =T_n-\frac1{2(2n+1)}>0,
\]
as is seen by replacing \(1/(1+x)\) by \(1/2\) in the integral.
Thus the coefficients \(q_n\) are nonnegative first and negative
thereafter, with at most one change of sign. They eventually are
negative, since the power average tends to zero whereas
\(2nT_n>n/(2n+1)\).

Let \(N\) be the last index with \(q_N\ge0\). Such an index exists
if the margin is nonnegative at \(r_1>0\). For
\(t_1=r_1^2\) and \(\lambda=t/t_1\in(0,1]\), the sign ordering gives
\[
 \sum_{n\ge1}q_nt^n\ge
             \lambda^N\sum_{n\ge1}q_nt_1^n\ge0.
\]
Indeed, \(\lambda^n\ge\lambda^N\) before the sign change; afterwards
the inequality reverses when multiplied by the negative coefficient.
This proves the propagation, with \(r=0\) immediate.

Finally \(\partial_a^2M_{\rm loc}(r,a)=-(1-r^2)r^2/(1-a^2r^2)<0\) and
\(M_{\rm loc}(r,1)=0\). A successful check at \((r_1,a_0)\) therefore covers
all \(r\le r_1\) and \(a\ge a_0\): concavity extends in \(a\),
and the series argument extends in \(r\).
